\documentclass[11pt]{article}
\usepackage[margin=1in]{geometry}
\usepackage[T1]{fontenc}
\usepackage{lmodern}
\usepackage{amsmath,amssymb,amsthm,mathtools}
\usepackage{booktabs,array}
\usepackage{microtype}
\usepackage{hyperref}
\hypersetup{colorlinks=true,linkcolor=blue,citecolor=blue,urlcolor=blue,
 pdftitle={A sum-of-squares proof of the four-qubit Higuchi--Sudbery entropy bound},
 pdfauthor={Dev D. Goyal}}
\setlength{\parskip}{4pt}
\setlength{\parindent}{0pt}

\newtheorem{theorem}{Theorem}[section]
\newtheorem{lemma}[theorem]{Lemma}
\newtheorem{proposition}[theorem]{Proposition}
\newtheorem{corollary}[theorem]{Corollary}
\theoremstyle{definition}
\newtheorem{definition}[theorem]{Definition}
\theoremstyle{remark}
\newtheorem{remark}[theorem]{Remark}
\DeclareMathOperator{\Tr}{Tr}
\newcommand{\C}{\mathbb C}
\newcommand{\R}{\mathbb R}
\newcommand{\N}{\mathbb N}
\newcommand{\HS}{\mathrm{HS}}
\newcommand{\ket}[1]{\lvert #1\rangle}
\newcommand{\conj}[1]{\overline{#1}}

\title{A sum-of-squares proof of the four-qubit\newline
Higuchi--Sudbery entropy bound}
\author{Dev D. Goyal}
\date{27 October 2026}

\begin{document}
\maketitle

\par\noindent{\normalsize\bfseries This work is fully autonomous agentic research. The initial research and proofs were generated autonomously by GPT 6 with Pro Reasoning in ChatGPT. GPT 6 Astra in Codex autonomously produced the formalizations and corrections. The authors checked the proofs to the best of their mathematical abilities. This check is not expert validation or peer review. Dev D. Goyal is responsible for publication.\par}

\begin{abstract}
For a pure state of four qubits, the Higuchi--Sudbery problem asks for the largest average von Neumann entropy across the three partitions into pairs. The proof shows that this average is at most $\log2+\tfrac12\log3$, with natural logarithms. The Higuchi--Sudbery state attains the bound. The proof gives an exact polynomial identity that bounds the sum of the third elementary symmetric functions of the three marginal spectra by $5/36$. Its right side is a sum of 272 squared mixed cubics with positive integer weights. The identity holds for all sixteen complex amplitudes. A purity lower bound and a cubic majorant for $-x\log x$ then give the entropy inequality. The argument places no restriction on the one-qubit marginals. A companion Lean formalization includes the polynomial identity, the actual reduced-matrix eigenvalues, and an explicit state that attains the bound.
\end{abstract}

\section{The problem and the contribution}

For a normalized pure state of four distinguishable qubits, let
$\rho_{AB}$, $\rho_{AC}$, and $\rho_{AD}$ denote the reduced density
matrices on the indicated pairs. The notation uses natural logarithms, with
\[
 S(\rho)=-\sum_i\lambda_i\log\lambda_i,
 \qquad 0\log0=0,
\]
where $\lambda_i$ are the eigenvalues of $\rho$. The quantity to be
maximized is
\begin{equation}
 \mathcal S(\psi)=\frac13\bigl(S(\rho_{AB})+S(\rho_{AC})+S(\rho_{AD})\bigr).
 \label{eq:objective}
\end{equation}

Higuchi and Sudbery introduced the candidate maximizing state in 2000 and
gave evidence for global maximality \cite{HS00}. Brierley and Higuchi
proved that it is a local maximum \cite{BH07}. Gour and Wallach proved the
corresponding global entropy statement on their family $\mathcal A$
\cite[Theorem~10(a)]{GW10}; by their Theorem~5, that family, up to local
unitaries, consists of states whose one-qubit marginals are all maximally
mixed. The hypothesis $\psi\in\mathcal A$ is material. The unrestricted
von Neumann entropy problem is not the statement of that theorem.
Recent primary sources continue to distinguish the restricted theorem
from the conjectured unrestricted optimum: see Pauwels and G\"uhne
\cite[Section~IV.2 and Appendix~A]{PG26} and Collins
\cite[Section~3.2]{C26}.

\begin{theorem}[Balanced four-qubit entropy bound]
\label{thm:main}
For every normalized $\psi\in(\C^2)^{\otimes4}$,
\begin{equation}
 \mathcal S(\psi)\leq\log2+\frac12\log3.
 \label{eq:main}
\end{equation}
Equality is attained by the Higuchi--Sudbery state in
Section~\ref{sec:sharpness}.
\end{theorem}

In base-two units the value is $1+\tfrac12\log_2 3$, approximately
$1.79248$ bits. The proof uses the unrestricted degree-six inequality in Theorem~\ref{thm:e3}, together with
its explicit certificate. The purity lower bound is known. A direct algebraic proof makes its normalization and use explicit. The maximizing state, its marginal spectra, and the restricted result are due to prior work.

\section{Coordinates and reduced matrices}

Write the unnormalized state as
\[
 z=\sum_{a,b,c,d\in\{0,1\}}z_{abcd}\ket{abcd},
 \qquad
 n=\|z\|^2=\sum_{a,b,c,d}|z_{abcd}|^2.
\]
For single-index formulas, $z_{abcd}$ means $z_{8a+4b+2c+d}$.
Define three flattenings, with rows and columns ordered lexicographically,
by
\begin{equation}
 (M_0)_{ab,cd}=z_{abcd},\qquad
 (M_1)_{ac,bd}=z_{abcd},\qquad
 (M_2)_{ad,bc}=z_{abcd}.
 \label{eq:flatten}
\end{equation}
Let $\rho_j=M_jM_j^*$. These matrices are positive semidefinite and have
trace $n$. When $n=1$, they are the three density matrices in
\eqref{eq:objective}. For each of these reduced matrices $A$, with
eigenvalues $\lambda_1,\ldots,\lambda_4$, put
\[
 e_3(A)=\sum_{1\leq i<j<k\leq4}\lambda_i\lambda_j\lambda_k.
\]
Equivalently, $e_3(A)$ is the sum of the four principal minors of order
three.

\begin{lemma}[Cofactor expression]
\label{lem:cofactor}
For every state vector $z$ and each flattening $M=M_j$ in
\eqref{eq:flatten},
\begin{equation}
 e_3(MM^*)=\sum_{\substack{I,J\subseteq\{0,1,2,3\}\\|I|=|J|=3}}
       |\det M[I,J]|^2.
 \label{eq:cofactor}
\end{equation}
\end{lemma}
\begin{proof}
Sum the principal minors $\det((MM^*)[I,I])$ over the four three-element
row sets $I$. For a fixed $I$, Cauchy--Binet expands this determinant as
$\sum_{|J|=3}\det M[I,J]\det M^*[J,I]$. The second factor is the
complex conjugate of the first, giving \eqref{eq:cofactor}.
\end{proof}

\section{An exact degree-six certificate}
\label{sec:sos}

\begin{theorem}[Third elementary symmetric function bound]
\label{thm:e3}
For every $z\in(\C^2)^{\otimes4}$, without a normalization assumption,
\begin{equation}
 \sum_{j=0}^2 e_3(M_jM_j^*)\leq\frac5{36}\|z\|^6.
 \label{eq:e3}
\end{equation}
The constant is sharp.
\end{theorem}

The proof consists of the following finite identity. The indexing and coefficients below specify the identity exactly.

\subsection{Local integer rows}

Each row in Table~\ref{tab:rows} is an integer-valued function on
$\{0,1\}^3$. A triple such as $001$ denotes its indicator function;
all omitted values are zero.

\begin{table}[ht]
\centering
\begin{tabular}{@{}cl@{}}
\toprule
Row & Nonzero values, expressed as indicators\\
\midrule
$S_0$ & $000$\\
$S_1$ & $001+010+100$\\
$S_2$ & $011+101+110$\\
$S_3$ & $111$\\
$T_0$ & $2\cdot001-010-100$\\
$T_1$ & $011+101-2\cdot110$\\
$A_0$ & $010-100$\\
$A_1$ & $011-101$\\
\bottomrule
\end{tabular}
\caption{Integer rows defining the squared cubics.}
\label{tab:rows}
\end{table}

Given a four-tuple of rows $r=(r_0,r_1,r_2,r_3)$, define
\begin{equation}
 C_r(z)=\sum_{a,b,c\in\{0,1\}^4}
       \left(\prod_{j=0}^3 r_j(a_j,b_j,c_j)\right)
       (-1)^{|c|}z_a z_b\conj{z_{\bar c}}.
 \label{eq:cubic}
\end{equation}
Here $|c|=c_0+c_1+c_2+c_3$, and $\bar c$ is the bitwise complement.
Thus $C_r$ has degree two in $z$ and degree one in $\bar z$.
The sign and complement in \eqref{eq:cubic} are exactly those of
$J^{\otimes4}\bar z$, where
$J=\left(\begin{smallmatrix}0&1\\-1&0\end{smallmatrix}\right)$.

Fix a site $s\in\{0,1,2,3\}$, an index $r\in\{0,1,2,3\}$, and a bit
$\epsilon_j\in\{0,1\}$ for each $j\ne s$. Define
$P_{s,r,\epsilon}$ by taking row $S_r$ at site $s$ and row
$T_{\epsilon_j}$ at each other site in \eqref{eq:cubic}. Define
$Q_{s,r,\epsilon}$ as the sum of the three cubics obtained from this
choice by replacing exactly two of its three $T$ rows with the
corresponding $A$ rows. Finally, for $\epsilon\in\{0,1\}^4$, let
$W_\epsilon$ be \eqref{eq:cubic} with row $A_{\epsilon_j}$ at every
site. Put
\begin{equation}
 w_0=w_3=3,\qquad w_1=w_2=1.
 \label{eq:weights}
\end{equation}

The reconstruction recipe first forms the polynomial expression
\begin{align*}
 R_{\mathrm{recipe}}(z)={}&
 \sum_{s=0}^3\sum_{r=0}^3\sum_{\epsilon\in\{0,1\}^{\{0,1,2,3\}\setminus\{s\}}}
 w_r\left(|5P_{s,r,\epsilon}+11Q_{s,r,\epsilon}|^2
             +124|Q_{s,r,\epsilon}|^2\right)
 \\
 &+640\sum_{\epsilon\in\{0,1\}^4}|W_\epsilon|^2.
\end{align*}
The generator divides by $16$, expands each cubic, removes its integer
common factor, and absorbs the square of that factor into its weight.
The resulting finite list is named \texttt{sosData} in
\texttt{HS/CertificateData.lean}, which assembles the explicit coefficient
lists in \path{HS/CertificateParts/L00.lean} through
\texttt{L16.lean}. It has 272 entries
$(v_\ell,T_\ell)$, where $v_\ell\in\{1,3,31,93,160\}$ and
$T_\ell$ is a list of integer-index tuples $(a,i,j,k)$.
Define directly from these delivered lists
\[
 F_\ell(z)=\sum_{(a,i,j,k)\in T_\ell}a z_i z_j\conj{z_k}.
\]
The row recipe permits an independent reconstruction of the data. The
formal theorem takes the explicit coefficient lists as its definitions
and checks the resulting identity directly. It does not verify the
Python generator or rely on a formal proof of the recipe-to-table
conversion.

\begin{samepage}
\begin{proposition}[Explicit-table polynomial identity]
\label{prop:certificate}
For all $z\in\C^{16}$, the delivered coefficient lists satisfy
\begin{equation}
 400n^3-2880
   \sum_{j=0}^2\sum_{\substack{|I|=|J|=3}}|\det M_j[I,J]|^2
 =\sum_{\ell=1}^{272}v_\ell|F_\ell(z)|^2.
 \label{eq:certificate}
\end{equation}
\end{proposition}
\end{samepage}

\begin{proof}
This is an identity of polynomials of bidegree $(3,3)$ in the amplitudes
and their conjugates. It can be checked by treating the sixteen amplitudes
and sixteen formal conjugates as independent commuting variables.
For precision, use the monomial basis
\[
 z_{i_1}z_{i_2}z_{i_3}\bar z_{j_1}\bar z_{j_2}\bar z_{j_3},
 \qquad i_1\leq i_2\leq i_3,\quad j_1\leq j_2\leq j_3.
\]
Expand $n^3$, expand each determinant by its six permutation terms,
and expand each cubic from its explicit list $T_\ell$.
Multiplication and collection of coefficients give
identical integer coefficients on the two sides of
\eqref{eq:certificate}. Each side has exactly 1,968 nonzero collected
coefficients; the difference has none.

The accompanying files supply all 272 cubics and all 48
determinants explicitly as coefficient lists. A separately implemented
checker forms $M_jM_j^*$ and expands its principal order-three minors
directly; it obtains the same polynomial before comparison with the
squared cubics. Thus the identity is specified by finite formulas here
and is independently reproducible without the optimization procedure
used to find it. The companion Lean files check the polynomial equalities
in the proof kernel. Section~\ref{sec:verification} describes the pinned
environment and verification records.

The Lean development specifies the integer coefficient lists explicitly.
Its kernel checks the resulting polynomial
identity and its nonnegativity consequence; it does not assume that the
Python generator, the discovery solver, or a representation-theoretic
decomposition is correct.
\end{proof}

\begin{proof}[Proof of Theorem~\ref{thm:e3}]
Every summand on the right of \eqref{eq:certificate} is nonnegative.
Divide by $2880$ and use $400/2880=5/36$, followed by
Lemma~\ref{lem:cofactor}. Sharpness follows from the explicit marginal
spectrum in Section~\ref{sec:sharpness}: its $e_3$ value is $5/108$
at each of the three cuts.
\end{proof}

\subsection{Equivalent moment form}

For a normalized state, set
\begin{equation}
 q_k=\frac13\sum_{j=0}^2\Tr(\rho_j^k),\qquad k=2,3.
 \label{eq:qk}
\end{equation}
Newton's identity, with $\Tr\rho_j=1$, reads
\begin{equation}
 6e_3(\rho_j)=1-3\Tr(\rho_j^2)+2\Tr(\rho_j^3).
 \label{eq:newton}
\end{equation}
Summing and applying \eqref{eq:e3} yields
\begin{equation}
 q_3\leq\frac32q_2-\frac{13}{36}.
 \label{eq:moment}
\end{equation}
For an unnormalized vector the equivalent inequality is
\[
 12\sum_j\Tr(\rho_j^3)+13n^3
 \leq18n\sum_j\Tr(\rho_j^2).
\]

\section{The purity lower bound}

The argument also uses $q_2\geq1/3$, which is a known four-qubit
purity bound. Its relation with the four-tangle appears explicitly in
\cite[Theorem~6]{GW10}. The following coordinate identity provides a
separate finite-algebra proof and fixes all normalization factors.

For each site $s$, let $z^{(s)}_0,z^{(s)}_1\in\C^8$ be the zero and one
slices of the amplitude tensor, in the same order of the remaining bits.
Define
\[
 d_s=\|z^{(s)}_0\|^2-\|z^{(s)}_1\|^2,
 \qquad c_s=\sum_u (z^{(s)}_0)_u\conj{(z^{(s)}_1)_u},
\]
and
\begin{equation}
 H=z_0z_{15}-z_1z_{14}-z_2z_{13}+z_3z_{12}
       -z_4z_{11}+z_5z_{10}+z_6z_9-z_7z_8.
 \label{eq:H}
\end{equation}

\begin{lemma}[Quartic purity identity]
\label{lem:purity}
For every unnormalized four-qubit vector,
\begin{equation}
 \sum_{j=0}^2\Tr(\rho_j^2)-n^2
 =\frac12\sum_{s=0}^3(d_s^2+4|c_s|^2)+2|H|^2.
 \label{eq:purityidentity}
\end{equation}
Consequently $q_2\geq1/3$ when $n=1$.
\end{lemma}
\begin{proof}
Substitute the flattenings \eqref{eq:flatten} in the left side and the
slice definitions and \eqref{eq:H} in the right side. Expansion gives the
same quartic polynomial. The identity is also valid with independent
formal conjugate variables, and this is the version used for the direct
ring calculation in the companion formalization. Every term on the right
of \eqref{eq:purityidentity} is nonnegative. Therefore the three purities
sum to at least $n^2$; divide by three and set $n=1$.
\end{proof}

\section{An entropy bound from two moments}

The remaining argument is a scalar interpolation calculation. The proof needs only the one-sided fourth-derivative statement below. Its hypotheses
and conclusion agree with the corresponding Lean lemma.

\begin{lemma}[A fourfold Rolle argument]
\label{lem:rolle}
Let $f_0,f_1,f_2,f_3,f_4:(0,\infty)\to\R$ satisfy
$f_j'(t)=f_{j+1}(t)$ for every $t>0$ and $j=0,1,2,3$, and suppose
$f_4(t)>0$ for every $t>0$. If $0<a<b$ and
\[
 f_0(a)=f_1(a)=f_0(b)=f_1(b)=0,
\]
then $f_0(x)\geq0$ for every $x>0$.
\end{lemma}
\begin{proof}
The conclusion at $a,b$ is immediate. Suppose, for a contradiction, that
$x>0$ is distinct from $a,b$ and $f_0(x)<0$. Put
\[
 q(t)=(t-a)^2(t-b)^2,\qquad k=\frac{f_0(x)}{q(x)}<0,
 \qquad g_j(t)=f_j(t)-kq^{(j)}(t)\quad(0\leq j\leq4).
\]
Then $g_j'=g_{j+1}$ for $j=0,1,2,3$; moreover, $g_0$ vanishes at
$a,b,x$ and $g_1$ vanishes at $a,b$. The argument shows that $g_1$ has four
\emph{distinct} positive zeros. If $x<a$, Rolle's theorem applied on
$[x,a]$ and $[a,b]$ supplies zeros $u\in(x,a)$ and $v\in(a,b)$;
the four zeros are $u,a,v,b$. If $a<x<b$, the intervals $[a,x]$ and
$[x,b]$ instead give four zeros $a,u,v,b$. If $b<x$, the intervals
$[a,b]$ and $[b,x]$ give four zeros $a,u,b,v$.

Three further rounds of Rolle's theorem, first between consecutive
zeros of $g_1$, then between the resulting zeros of $g_2$, and finally
between the resulting zeros of $g_3$, produce a positive zero of $g_4$.
All intervals used lie inside $(0,\infty)$; differentiability supplies
the continuity needed by Rolle's theorem on each closed interval. This
is impossible because
\[
 g_4(t)=f_4(t)-24k>0\qquad(t>0).
\]
\end{proof}

Let $L=\log3$ and define
\begin{align}
 \alpha&=\frac12-\frac38L,&
 \beta&=\log2+6L-\frac{11}{2},\notag\\
 \gamma&=18-\frac{39}{2}L,&
 \delta&=18(L-1),
 \label{eq:coefficients}
\end{align}
with $p(x)=\alpha+\beta x+\gamma x^2+\delta x^3$.

The required elementary logarithmic bounds are
\begin{equation}
 1<\log3<\frac65.
 \label{eq:logbounds}
\end{equation}
To check them using exactly the rational comparisons in the Lean proof,
use $1-t^{-1}\leq\log t\leq t-1$ for $t>0$ together with
\[
 \left(\frac{12}{11}\right)^{12}<3
 <\left(\frac{23}{20}\right)^8.
\]
The first comparison and monotonicity of the logarithm give
\[
 1=12\left(1-\frac{11}{12}\right)
 \leq12\log\frac{12}{11}<\log3.
\]
The second gives
\[
 \log3<8\log\frac{23}{20}
 \leq8\left(\frac{23}{20}-1\right)=\frac65.
\]
In particular,
\begin{equation}
 \alpha>0,\qquad\delta>0,\qquad
 \gamma+\frac32\delta=\frac{15}{2}\log3-9<0.
 \label{eq:signs}
\end{equation}

\begin{lemma}[Scalar entropy majorant]
\label{lem:majorant}
For all $x\geq0$, with $0\log0=0$,
\begin{equation}
 -x\log x\leq p(x).
 \label{eq:majorant}
\end{equation}
The two sides agree at $x=1/6$ and $x=1/2$.
\end{lemma}
\begin{proof}
For $x>0$, take $f_0(x)=p(x)+x\log x$, with derivative chain
\begin{align*}
 f_1(x)&=\beta+2\gamma x+3\delta x^2+\log x+1,\\
 f_2(x)&=2\gamma+6\delta x+\frac1x,\\
 f_3(x)&=6\delta-\frac1{x^2},\\
 f_4(x)&=\frac2{x^3}>0.
\end{align*}
Substitution of \eqref{eq:coefficients}, together with
$\log(1/6)=-\log2-\log3$ and $\log(1/2)=-\log2$, gives
\[
 f_0(1/6)=f_1(1/6)=f_0(1/2)=f_1(1/2)=0.
\]
Lemma~\ref{lem:rolle} therefore gives $f_0(x)\geq0$ for all $x>0$,
which is \eqref{eq:majorant}. The same substitutions prove equality at
the two specified points. At $x=0$ the left-hand side is zero and
$p(0)=\alpha>0$ by \eqref{eq:logbounds}, so the inequality holds there
as well.
\end{proof}

\begin{proof}[Proof of the upper bound in Theorem~\ref{thm:main}]
Apply \eqref{eq:majorant} to each of the four nonnegative eigenvalues
of each density matrix. Their sum is one, so
\[
 S(\rho_j)\leq4\alpha+\beta+
       \gamma\Tr(\rho_j^2)+\delta\Tr(\rho_j^3).
\]
Averaging and using \eqref{eq:moment}, with $\delta>0$, yields
\[
 \mathcal S(\psi)\leq4\alpha+\beta+
       (\gamma+\tfrac32\delta)q_2-\tfrac{13}{36}\delta.
\]
The coefficient of $q_2$ is negative by \eqref{eq:signs}; hence
Lemma~\ref{lem:purity} gives
\begin{align*}
 \mathcal S(\psi)
 &\leq4\alpha+\beta+
       \tfrac13(\gamma+\tfrac32\delta)-\tfrac{13}{36}\delta\\
 &=\log2+\tfrac12\log3,
\end{align*}
where the last equality is substitution of \eqref{eq:coefficients}.
\end{proof}

\section{Normalization and attainment}
\label{sec:sharpness}

Let $\omega=(-1+i\sqrt3)/2$ and set
\begin{equation}
 \ket{\HS}=\frac1{\sqrt6}\left[
 \ket{0011}+\ket{1100}
 +\omega(\ket{0101}+\ket{1010})
 +\omega^2(\ket{0110}+\ket{1001})\right].
 \label{eq:state}
\end{equation}
The six basis vectors are distinct and the six coefficients have modulus
$1/\sqrt6$, so the state is normalized. For exact real arithmetic the
same vector has two amplitudes equal to $2$, two equal to
$-1+i\sqrt3$, and two equal to $-1-i\sqrt3$, all divided by $\sqrt{24}$.
This is the representation used in the formal witness.

For the cut $AB|CD$, the unnormalized flattening of the bracketed
expression in \eqref{eq:state} is
\[
 \begin{pmatrix}
 0&0&0&1\\
 0&\omega&\omega^2&0\\
 0&\omega^2&\omega&0\\
 1&0&0&0
 \end{pmatrix}.
\]
The other two cuts replace the isolated unit amplitudes and the middle
pair by the corresponding phases. In all three products with the
conjugate transpose, using $\omega+\omega^2=-1$ gives
\begin{equation}
 \rho_0=\rho_1=\rho_2=
 \frac16\begin{pmatrix}
 1&0&0&0\\
 0&2&-1&0\\
 0&-1&2&0\\
 0&0&0&1
 \end{pmatrix}.
 \label{eq:HSdensity}
\end{equation}
The middle block has eigenvectors $(1,1)$ and $(1,-1)$ with
eigenvalues $1$ and $3$, respectively. Thus the eigenvalues of each
density matrix are $1/2,1/6,1/6,1/6$, and
\[
 S(\rho_j)=-\tfrac12\log(\tfrac12)-\tfrac12\log(\tfrac16)
 =\log2+\tfrac12\log3.
\]
This proves attainment in Theorem~\ref{thm:main}. The same spectrum gives
$e_3(\rho_j)=5/108$, proving sharpness in Theorem~\ref{thm:e3}.
No classification of all maximizing vectors is asserted here.

\section{Formalization and scope}
\label{sec:verification}

The companion formalization uses Lean~4.35.0 and a pinned mathlib revision. Its final greatest-value statements include both the upper bounds and their attainment. They use the eigenvalues of the actual reduced matrices. The reproduction package gives the exact source versions, build commands, and axiom checks.

The coefficient data can also be checked with integer arithmetic. One program reconstructs the cubics from the row formulas and expands the forty-eight determinants. A separate checker directly expands the principal minors of the reduced density matrices. Both compare the collected coefficients. This check does not require a numerical optimization solver.

The result applies to pure states of four distinguishable two-dimensional systems. It uses the equally weighted average of the three balanced cuts. Other local dimensions, mixed states, other entanglement measures, and different cut weights require separate results. The proof gives an attaining state but does not classify all maximizers.


\begin{thebibliography}{9}

\bibitem{HS00}
A.~Higuchi and A.~Sudbery,
\emph{How entangled can two couples get?},
Physics Letters A \textbf{273} (2000), 213--217.
\href{https://doi.org/10.1016/S0375-9601(00)00480-1}{doi:10.1016/S0375-9601(00)00480-1}.
\url{https://arxiv.org/abs/quant-ph/0005013}.

\bibitem{BH07}
S.~Brierley and A.~Higuchi,
\emph{On maximal entanglement between two pairs in four-qubit pure states},
Journal of Physics A: Mathematical and Theoretical \textbf{40} (2007),
8455--8465.
\href{https://doi.org/10.1088/1751-8113/40/29/018}{doi:10.1088/1751-8113/40/29/018}.
\url{https://arxiv.org/abs/0704.2961}.

\bibitem{GW10}
G.~Gour and N.~R.~Wallach,
\emph{All maximally entangled four-qubit states},
Journal of Mathematical Physics \textbf{51} (2010), 112201.
\href{https://doi.org/10.1063/1.3511477}{doi:10.1063/1.3511477}.
\url{https://arxiv.org/abs/1006.0036}.

\bibitem{PG26}
J.~Pauwels and O.~G\"uhne,
\emph{The Entanglement Content of Quantum Measurement Bases},
arXiv preprint (2026).
\url{https://arxiv.org/abs/2608.21185}.

\bibitem{C26}
B.~Collins,
\emph{The Failure of Simultaneous Quantum Typicality},
arXiv preprint (2026).
\url{https://arxiv.org/abs/2609.25726}.

\end{thebibliography}
\end{document}
